% ============================================================
% 培优专题16 — 组合计数进阶
% 编译命令：xelatex -interaction=nonstopmode 培优专题16_组合计数进阶.tex
% ============================================================
\documentclass[12pt, a4paper]{ctexart}
\usepackage[margin=2.2cm]{geometry}
\usepackage{amsmath, amssymb, amsfonts}
\usepackage{enumitem}
\usepackage{tcolorbox}
\usepackage{booktabs}
\usepackage{array}
\usepackage{fancyhdr}
\usepackage{xcolor}
\usepackage{tikz}
\usepackage{hyperref}
\usepackage{graphicx}
\usepackage{multirow}
\usepackage{pgfplots}
\pgfplotsset{compat=1.18}

% 颜色定义
\definecolor{mainblue}{RGB}{41, 98, 168}
\definecolor{lightblue}{RGB}{230, 240, 255}
\definecolor{tipgreen}{RGB}{34, 139, 87}
\definecolor{lightgreen}{RGB}{235, 250, 240}
\definecolor{warnred}{RGB}{200, 50, 50}
\definecolor{lightred}{RGB}{255, 240, 240}
\definecolor{examorange}{RGB}{210, 120, 20}
\definecolor{lightorange}{RGB}{255, 248, 235}
\definecolor{apppurple}{RGB}{120, 60, 170}
\definecolor{lightpurple}{RGB}{245, 238, 255}

% 页眉页脚
\pagestyle{fancy}
\fancyhf{}
\fancyhead[L]{\small\textcolor{mainblue}{\textbf{高中数学培优 \(\cdot\) 组合计数进阶}}}
\fancyhead[R]{\small\textcolor{mainblue}{培优专题16}}
\fancyfoot[C]{\small —\ \thepage\ —}
\renewcommand{\headrulewidth}{0.6pt}
\renewcommand{\headrule}{\hbox to\headwidth{\color{mainblue}\leaders\hrule height \headrulewidth\hfill}}
\renewcommand{\footrulewidth}{0pt}

% tcolorbox 样式
\newtcolorbox{knowledgebox}[1][]{%
  colback=lightblue, colframe=mainblue,
  fonttitle=\bfseries\large, title=#1,
  arc=3mm, boxrule=0.8pt,
  left=8pt, right=8pt, top=6pt, bottom=6pt
}

\newtcolorbox{tipbox}[1][]{%
  colback=lightgreen, colframe=tipgreen,
  fonttitle=\bfseries, title=#1,
  arc=2mm, boxrule=0.6pt,
  left=6pt, right=6pt, top=4pt, bottom=4pt
}

\newtcolorbox{warnbox}[1][]{%
  colback=lightred, colframe=warnred,
  fonttitle=\bfseries, title=#1,
  arc=2mm, boxrule=0.6pt,
  left=6pt, right=6pt, top=4pt, bottom=4pt
}

\newtcolorbox{examplebox}[1][]{%
  colback=lightorange, colframe=examorange,
  fonttitle=\bfseries, title=#1,
  arc=3mm, boxrule=0.8pt,
  left=8pt, right=8pt, top=6pt, bottom=6pt
}

\newtcolorbox{appbox}[1][]{%
  colback=lightpurple, colframe=apppurple,
  fonttitle=\bfseries, title=#1,
  arc=3mm, boxrule=0.8pt,
  left=8pt, right=8pt, top=6pt, bottom=6pt
}

\setCJKmainfont{Hiragino Sans GB}
\setCJKsansfont{Heiti SC}

\begin{document}

% ============ 标题页 ============
\begin{center}
{\Huge\bfseries\textcolor{mainblue}{培优专题16\quad 组合计数进阶}}\\[12pt]
{\Large\textcolor{gray}{高中数学培优系列}}\\[8pt]
\rule{0.6\textwidth}{1pt}\\[8pt]
{\normalsize 适用：高二高三拔高\textbackslash 竞赛入门 \(\quad\mid\quad\) 课时：3课时}
\end{center}
\vspace{10pt}

% ============================================================
% 第一部分：知识点讲解
% ============================================================
\section*{\textcolor{mainblue}{一、知识点讲解}}

% ========== 1.1 容斥原理 ==========
\subsection*{1.1\quad 容斥原理}

\begin{knowledgebox}[容斥原理（Inclusion-Exclusion Principle）]
设 \(A_1, A_2, \ldots, A_n\) 是有限集合，则
\[
\left|\bigcup_{i=1}^{n} A_i\right|
= \sum_{i} |A_i| - \sum_{i<j} |A_i \cap A_j| + \sum_{i<j<k} |A_i \cap A_j \cap A_k| - \cdots + (-1)^{n+1}|A_1 \cap A_2 \cap \cdots \cap A_n|
\]

\textbf{常用特殊形式：}

\textbf{两集合：} \(|A \cup B| = |A| + |B| - |A \cap B|\)

\textbf{三集合：}
\[
|A \cup B \cup C| = |A|+|B|+|C| - |A\cap B|-|A\cap C|-|B\cap C| + |A\cap B\cap C|
\]

\textbf{补集形式（求"全不"的个数）：}
\[
\left|\overline{A_1} \cap \overline{A_2} \cap \cdots \cap \overline{A_n}\right|
= |U| - \left|\bigcup_{i=1}^{n} A_i\right|
\]
其中 \(U\) 为全集。
\end{knowledgebox}

\begin{warnbox}[易错提醒]
\begin{itemize}[leftmargin=1.5em]
  \item 容斥原理的关键是"不重不漏"：先加后减再加，交替进行。
  \item 计算交集大小时，要确认各约束条件是否独立（可乘则乘）。
  \item 应用补集形式时，"至少一个"用并集，"全都不"用补集。
\end{itemize}
\end{warnbox}

% ========== 1.2 递推计数 ==========
\subsection*{1.2\quad 递推计数}

\begin{knowledgebox}[递推计数方法]
\textbf{核心思想：} 将规模为 \(n\) 的问题分解为规模更小的子问题，建立递推关系。

\textbf{建立递推关系的一般步骤：}
\begin{enumerate}[leftmargin=1.5em]
  \item 确定状态：用 \(a_n\) 表示规模为 \(n\) 时的方案数。
  \item 分类讨论：根据最后一步（或某个关键决策）进行分类。
  \item 利用加法原理写出递推式。
  \item 确定初始条件（base case）。
\end{enumerate}

\textbf{常见递推模型：}
\begin{itemize}[leftmargin=1.5em]
  \item 斐波那契型：\(a_n = a_{n-1} + a_{n-2}\)，如铺砖问题、兔子繁殖。
  \item 线性递推：\(a_n = c_1 a_{n-1} + c_2 a_{n-2} + \cdots\)，如传球问题。
  \item 非线性递推：\(a_n = \sum_{k=0}^{n-1} a_k a_{n-1-k}\)，如卡特兰数。
\end{itemize}
\end{knowledgebox}

\begin{tipbox}[规律]
递推计数与分类计数的本质相同：递推是用"最后一步"分类，将全局计数化为子问题计数之和。学会画"决策树"辅助分析。
\end{tipbox}

% ========== 1.3 卡特兰数 ==========
\subsection*{1.3\quad 卡特兰数}

\begin{knowledgebox}[卡特兰数（Catalan Number）]
\textbf{定义：} 第 \(n\) 个卡特兰数为
\[
C_n = \frac{1}{n+1}\binom{2n}{n} = \frac{(2n)!}{(n+1)!\,n!}
\]

\textbf{递推公式：}
\[
C_0 = 1, \qquad C_n = \sum_{i=0}^{n-1} C_i \cdot C_{n-1-i} \quad (n \ge 1)
\]

\textbf{前几项：} \(C_0=1,\ C_1=1,\ C_2=2,\ C_3=5,\ C_4=14,\ C_5=42,\ C_6=132, \ldots\)

\textbf{典型模型（\(n\) 个卡特兰数 \(C_n\) 对应的计数问题）：}
\begin{enumerate}[leftmargin=1.5em]
  \item \(n\) 对括号的合法匹配方案数。
  \item 从 \((0,0)\) 到 \((n,n)\) 不越过对角线的格路数。
  \item \(n+1\) 个节点的不同构二叉树个数。
  \item 凸 \(n+2\) 边形的三角剖分数。
  \item 长度为 \(2n\) 的合法括号序列数。
  \item 出栈序列数：\(1,2,\ldots,n\) 依次入栈，可能的出栈序列数。
\end{enumerate}
\end{knowledgebox}

\begin{warnbox}[易错提醒]
\begin{itemize}[leftmargin=1.5em]
  \item 卡特兰数的下标从 \(C_0=1\) 开始，注意 \(n\) 的含义。
  \item 不同模型中"n"的含义不同：括号序列用 \(n\) 对，凸多边形用 \(n+2\) 条边。
  \item 卡特兰数的递推公式是"卷积"形式，不要与斐波那契递推混淆。
\end{itemize}
\end{warnbox}

% ========== 1.4 错排问题 ==========
\subsection*{1.4\quad 错排问题}

\begin{knowledgebox}[错排问题（Derangement）]
\textbf{定义：} \(n\) 个元素的\textbf{错排}是指一个排列 \(\sigma\)，使得对所有 \(i=1,2,\ldots,n\)，均有 \(\sigma(i) \ne i\)。即没有任何元素在原来的位置上。

\textbf{错排数公式（容斥原理推导）：}
\[
D_n = n!\sum_{k=0}^{n}\frac{(-1)^k}{k!} = n!\left(1 - \frac{1}{1!} + \frac{1}{2!} - \frac{1}{3!} + \cdots + \frac{(-1)^n}{n!}\right)
\]

\textbf{递推公式：}
\[
D_1 = 0,\quad D_2 = 1, \qquad D_n = (n-1)(D_{n-1} + D_{n-2}) \quad (n \ge 3)
\]

\textbf{前几项：} \(D_1=0,\ D_2=1,\ D_3=2,\ D_4=9,\ D_5=44,\ D_6=265, \ldots\)

\textbf{渐近公式：} 当 \(n\) 较大时，\(D_n \approx \dfrac{n!}{e}\)，即错排的概率趋近于 \(\dfrac{1}{e} \approx 0.3679\)。
\end{knowledgebox}

\begin{tipbox}[规律]
错排递推公式的直观理解：元素 \(n\) 可以放到位置 \(k\)（共 \(n-1\) 种选择）。若元素 \(k\) 恰好放到位置 \(n\)（互换），则剩余 \(n-2\) 个元素错排，有 \(D_{n-2}\) 种；若元素 \(k\) 不放到位置 \(n\)，则剩余 \(n-1\) 个元素（含元素 \(k\)）都不到各自位置，有 \(D_{n-1}\) 种。故 \(D_n = (n-1)(D_{n-1}+D_{n-2})\)。
\end{tipbox}

% ========== 1.5 路径计数 ==========
\subsection*{1.5\quad 路径计数}

\begin{knowledgebox}[格路计数]
\textbf{基本格路：} 从 \((0,0)\) 到 \((m,n)\)，每次向右或向上走一步，方案数为
\[
\binom{m+n}{m} = \binom{m+n}{n}
\]

\textbf{不越过对角线的格路（卡特兰数）：} 从 \((0,0)\) 到 \((n,n)\)，每次向右或向上走一步，且路径始终在对角线 \(y=x\) 下方（含对角线），方案数为卡特兰数 \(C_n = \dfrac{1}{n+1}\dbinom{2n}{n}\)。

\textbf{带障碍物的格路：} 利用容斥原理或递推（动态规划）计算。设 \(f(i,j)\) 为到达 \((i,j)\) 的方案数，则
\[
f(i,j) = f(i-1,j) + f(i,j-1)
\]
在障碍物处令 \(f(i,j)=0\)。

\textbf{反射原理（Andr\'{e} 反射）：} 从 \((0,0)\) 到 \((a,b)\)（其中 \(a>b\)）越过对角线的路径数等于从 \((-1,1)\) 到 \((a,b)\) 的路径数，即 \(\dbinom{a+b}{a+1}\)。
\end{knowledgebox}

\begin{warnbox}[易错提醒]
\begin{itemize}[leftmargin=1.5em]
  \item 格路计数中，"不越过对角线"与"不触碰对角线（除起点外）"是不同的条件，对应的计数结果不同。
  \item 反射原理仅适用于单一约束条件（不越过一条直线），多重约束需要递推或容斥。
\end{itemize}
\end{warnbox}

% ========== 1.6 生成函数入门 ==========
\subsection*{1.6\quad 生成函数入门}

\begin{knowledgebox}[普通生成函数（Ordinary Generating Function）]
\textbf{定义：} 数列 \(\{a_n\}_{n\ge 0}\) 的普通生成函数为
\[
A(x) = \sum_{n=0}^{\infty} a_n x^n = a_0 + a_1 x + a_2 x^2 + \cdots
\]

\textbf{常用生成函数：}
\begin{itemize}[leftmargin=1.5em]
  \item 常数数列 \(a_n=1\)：\(\displaystyle\frac{1}{1-x} = \sum_{n=0}^{\infty} x^n\)
  \item 等比数列 \(a_n=r^n\)：\(\displaystyle\frac{1}{1-rx} = \sum_{n=0}^{\infty} r^n x^n\)
  \item 组合数：\(\displaystyle\binom{m}{n}\) 的生成函数为 \((1+x)^m\)
  \item 卡特兰数：\(C(x) = \displaystyle\frac{1-\sqrt{1-4x}}{2x}\)
\end{itemize}

\textbf{生成函数的运算：}
\begin{itemize}[leftmargin=1.5em]
  \item \textbf{加法：} 对应数列逐项相加，即"或"的关系。
  \item \textbf{乘法（卷积）：} \(A(x) \cdot B(x)\) 对应数列的卷积 \(c_n = \sum_{k=0}^{n} a_k b_{n-k}\)，即"分步"的关系。
  \item \textbf{提取系数：} \([x^n] A(x)\) 表示 \(A(x)\) 中 \(x^n\) 的系数，即 \(a_n\)。
\end{itemize}
\end{knowledgebox}

\begin{tipbox}[规律]
生成函数的核心价值在于将离散的计数问题转化为代数运算。整数分拆、重复组合、递推求解等问题都可以用生成函数优雅地处理。关键是正确写出每个"零件"的生成函数，然后利用乘法（分步乘法原理）组合。
\end{tipbox}

% ============================================================
% 第二部分：例题精讲
% ============================================================
\section*{\textcolor{examorange}{二、例题精讲}}

% ----- 例题1：容斥原理基础 -----
\begin{examplebox}[【例题1】容斥原理——三集合基础]
某班有 50 名学生，其中选修数学的有 30 人，选修物理的有 25 人，选修化学的有 20 人；同时选修数学和物理的有 12 人，同时选修数学和化学的有 10 人，同时选修物理和化学的有 8 人；三门都选的有 5 人。求三门都不选的人数。
\end{examplebox}

\textbf{解：}

设 \(M, P, C\) 分别表示选修数学、物理、化学的学生集合。

由容斥原理（三集合形式）：
\begin{align*}
|M \cup P \cup C| &= |M|+|P|+|C| - |M\cap P|-|M\cap C|-|P\cap C| + |M\cap P\cap C| \\
&= 30+25+20 - 12-10-8 + 5 \\
&= 75 - 30 + 5 = 50
\end{align*}

三门都不选的人数为：
\[
|U| - |M \cup P \cup C| = 50 - 50 = \boxed{0}
\]

\begin{tipbox}[方法总结]
三集合容斥原理直接套公式，注意各项符号交替变化。验算技巧：代入极端值（如交集为零）检验公式正确性。
\end{tipbox}

\vspace{6pt}

% ----- 例题2：容斥原理——错位排列计数 -----
\begin{examplebox}[【例题2】容斥原理——满足多个条件的排列]
用 \(0,1,2,3,4,5\) 组成无重复数字的四位数，其中个位不是 \(1\)、十位不是 \(2\)、百位不是 \(3\) 的四位数有多少个？
\end{examplebox}

\textbf{解：}

设全集 \(U\) 为用 \(0,1,2,3,4,5\) 组成无重复数字的全部四位数。千位不能为 \(0\)，故
\[
|U| = 5 \times 5 \times 4 \times 3 = 300
\]

设 \(A\)：个位为 \(1\) 的四位数；\(B\)：十位为 \(2\) 的四位数；\(C\)：百位为 \(3\) 的四位数。

我们需求 \(|\overline{A} \cap \overline{B} \cap \overline{C}| = |U| - |A \cup B \cup C|\)。

\textbf{计算各集合大小：}

\(|A|\)（个位为 \(1\)）：千位从 \(\{2,3,4,5\}\) 中选（不能为 \(0\)），4 种；剩余两位从剩余 4 个数选 2 个排列，\(4 \times 3=12\) 种。故 \(|A|=4\times 12=48\)。

\(|B|\)（十位为 \(2\)）：千位从 \(\{1,3,4,5\}\) 中选（不能为 \(0\)），4 种；剩余两位从剩余 4 个数选 2 个排列，\(4\times 3=12\) 种。故 \(|B|=48\)。

\(|C|\)（百位为 \(3\)）：同理 \(|C|=48\)。

\(|A\cap B|\)（个位 \(1\)，十位 \(2\)）：千位从 \(\{3,4,5\}\) 中选，3 种；个位确定、十位确定，百位从剩余 3 个数选 1 个，3 种。故 \(|A\cap B|=3\times 3=9\)。

\(|A\cap C|\)（个位 \(1\)，百位 \(3\)）：千位从 \(\{2,4,5\}\) 中选，3 种；十位从剩余 3 个数选，3 种。故 \(|A\cap C|=9\)。

\(|B\cap C|\)（十位 \(2\)，百位 \(3\)）：千位从 \(\{1,4,5\}\) 中选，3 种；个位从剩余 3 个数选，3 种。故 \(|B\cap C|=9\)。

\(|A\cap B\cap C|\)（个位 \(1\)，十位 \(2\)，百位 \(3\)）：千位从 \(\{4,5\}\) 中选，2 种。故 \(|A\cap B\cap C|=2\)。

由容斥原理：
\begin{align*}
|A\cup B\cup C| &= 48+48+48 - 9-9-9 + 2 \\
&= 144 - 27 + 2 = 119
\end{align*}

故所求四位数个数为：
\[
|U| - |A\cup B\cup C| = 300 - 119 = \boxed{181}
\]

\begin{tipbox}[方法总结]
"某些位置不取某值"类问题，先计算"至少一个位置违反约束"的方案数（容斥原理），再从总数中减去。关键是正确计算各交集大小，注意首位不能为 \(0\) 的约束。
\end{tipbox}

\vspace{6pt}

% ----- 例题3：容斥原理——不定方程整数解 -----
\begin{examplebox}[【例题3】容斥原理——不定方程的约束解]
求方程 \(x_1 + x_2 + x_3 + x_4 = 20\) 满足 \(x_i \ge 0\) 且 \(x_1 \le 6, x_2 \le 8, x_3 \le 5\) 的整数解的个数。
\end{examplebox}

\textbf{解：}

设 \(U\) 为 \(x_1+x_2+x_3+x_4=20\)，\(x_i\ge 0\) 的全部非负整数解。由隔板法（插空法）：
\[
|U| = \binom{20+4-1}{4-1} = \binom{23}{3} = 1771
\]

设 \(A_1\)：\(x_1 \ge 7\)；\(A_2\)：\(x_2 \ge 9\)；\(A_3\)：\(x_3 \ge 6\)。

\(|A_1|\)：令 \(x_1' = x_1 - 7 \ge 0\)，则 \(x_1'+x_2+x_3+x_4 = 13\)，故 \(|A_1| = \dbinom{16}{3} = 560\)。

\(|A_2|\)：令 \(x_2' = x_2 - 9 \ge 0\)，则 \(x_1+x_2'+x_3+x_4 = 11\)，故 \(|A_2| = \dbinom{14}{3} = 364\)。

\(|A_3|\)：令 \(x_3' = x_3 - 6 \ge 0\)，则 \(x_1+x_2+x_3'+x_4 = 14\)，故 \(|A_3| = \dbinom{17}{3} = 680\)。

\(|A_1\cap A_2|\)：\(x_1'+x_2'+x_3+x_4 = 4\)，故 \(|A_1\cap A_2| = \dbinom{7}{3} = 35\)。

\(|A_1\cap A_3|\)：\(x_1'+x_2+x_3'+x_4 = 7\)，故 \(|A_1\cap A_3| = \dbinom{10}{3} = 120\)。

\(|A_2\cap A_3|\)：\(x_1+x_2'+x_3'+x_4 = 5\)，故 \(|A_2\cap A_3| = \dbinom{8}{3} = 56\)。

\(|A_1\cap A_2\cap A_3|\)：\(x_1'+x_2'+x_3'+x_4 = -2\)，无解，故 \(= 0\)。

由容斥原理，所求解数为：
\begin{align*}
&\quad |U| - |A_1\cup A_2\cup A_3| \\
&= 1771 - (560+364+680) + (35+120+56) - 0 \\
&= 1771 - 1604 + 211 = \boxed{378}
\end{align*}

\begin{tipbox}[方法总结]
不定方程约束解的标准做法：先算无约束的总数（隔板法），再用容斥原理减去违反上界约束的部分。"违反上界 \(x_i \le m\)"即"\(x_i \ge m+1\)"，做变量替换 \(x_i' = x_i - (m+1)\) 后转化为标准隔板法。
\end{tipbox}

\vspace{6pt}

% ----- 例题4：容斥原理——欧拉函数 -----
\begin{examplebox}[【例题4】容斥原理——欧拉函数]
利用容斥原理求 \(\varphi(360)\)，即不超过 \(360\) 且与 \(360\) 互质的正整数的个数。
\end{examplebox}

\textbf{解：}

首先分解质因数：\(360 = 2^3 \times 3^2 \times 5\)。

与 \(360\) 不互质等价于能被 \(2\)、\(3\) 或 \(5\) 中至少一个整除。设 \(A\)：被 \(2\) 整除的集合；\(B\)：被 \(3\) 整除的集合；\(C\)：被 \(5\) 整除的集合。

\begin{align*}
|A| &= \lfloor 360/2 \rfloor = 180, \quad |B| = \lfloor 360/3 \rfloor = 120, \quad |C| = \lfloor 360/5 \rfloor = 72 \\
|A\cap B| &= \lfloor 360/6 \rfloor = 60, \quad |A\cap C| = \lfloor 360/10 \rfloor = 36, \quad |B\cap C| = \lfloor 360/15 \rfloor = 24 \\
|A\cap B\cap C| &= \lfloor 360/30 \rfloor = 12
\end{align*}

不与 \(360\) 互质的个数：
\[
|A\cup B\cup C| = 180+120+72 - 60-36-24 + 12 = 264
\]

因此
\[
\varphi(360) = 360 - 264 = \boxed{96}
\]

\textbf{验证（欧拉函数公式）：}
\[
\varphi(360) = 360\left(1-\frac{1}{2}\right)\left(1-\frac{1}{3}\right)\left(1-\frac{1}{5}\right) = 360 \times \frac{1}{2} \times \frac{2}{3} \times \frac{4}{5} = 96 \quad \checkmark
\]

\begin{tipbox}[方法总结]
欧拉函数 \(\varphi(n)\) 本质上是容斥原理的应用。公式 \(\varphi(n) = n\prod_{p|n}\left(1-\frac{1}{p}\right)\) 可直接由容斥原理推导。质因数分解是前提。
\end{tipbox}

\vspace{6pt}

% ----- 例题5：递推计数——铺砖问题 -----
\begin{examplebox}[【例题5】递推计数——铺砖问题]
用 \(1\times 1\) 的方砖和 \(1\times 2\) 的长砖铺满一个 \(1\times n\) 的长条区域，有多少种铺法？
\end{examplebox}

\textbf{解：}

设铺满 \(1\times n\) 区域的铺法数为 \(a_n\)。

\textbf{建立递推：} 考虑最右端的铺法：
\begin{itemize}[leftmargin=1.5em]
  \item 若最右端放一块 \(1\times 1\) 方砖，则左边有 \(a_{n-1}\) 种铺法。
  \item 若最右端放一块 \(1\times 2\) 长砖（横放），则左边有 \(a_{n-2}\) 种铺法。
\end{itemize}

因此 \(a_n = a_{n-1} + a_{n-2}\)，这正是\textbf{斐波那契递推}。

\textbf{初始条件：} \(a_1 = 1\)（只放一块方砖），\(a_2 = 2\)（两块方砖或一块长砖）。

\textbf{计算：}
\[
a_1=1,\ a_2=2,\ a_3=3,\ a_4=5,\ a_5=8,\ a_6=13,\ \ldots
\]

即 \(a_n = F_{n+1}\)（第 \(n+1\) 个斐波那契数），其中 \(F_1=F_2=1\)。

铺满 \(1\times n\) 的铺法数为 \(\boxed{F_{n+1}}\)，即第 \(n+1\) 个斐波那契数。

\begin{tipbox}[方法总结]
铺砖问题的本质是"最后一步"分类：最右端放什么决定了剩余部分的规模。得到斐波那契递推后，可尝试求通项公式（特征方程法），也可直接用递推计算具体值。
\end{tipbox}

\vspace{6pt}

% ----- 例题6：递推计数——传球问题 -----
\begin{examplebox}[【例题6】递推计数——传球问题]
\(n\) 个人围成一圈传球，从某人开始，每次可以将球传给左边或右边相邻的人。经过 \(k\) 次传球后，球回到出发者的传球方式有多少种？
\end{examplebox}

\textbf{解：}

设 \(a_k\) 表示经 \(k\) 次传球后球回到出发者的方式数，\(b_k\) 表示经 \(k\) 次传球后球在某个特定非出发者的方式数（由对称性，球在任意一个非出发者处的方式数相同）。

\textbf{建立递推：}
\begin{itemize}[leftmargin=1.5em]
  \item 球在出发者时，上一步必须在左边或右边的邻居处：\(a_k = 2b_{k-1}\)。
  \item 球在某非出发者时，上一步可从该非出发者的两个邻居来：一个邻居是出发者（有 \(a_{k-1}\) 种），另一个是非出发者（有 \(b_{k-1}\) 种）。故 \(b_k = a_{k-1} + b_{k-1}\)。
\end{itemize}

初始条件：\(a_0 = 1\)（起始就在出发者处），\(b_0 = 0\)。

对于 \textbf{\(n=3\)} 的特殊情况（三人传球）：

由 \(a_k = 2b_{k-1}\) 和 \(b_k = a_{k-1}+b_{k-1}\)，代入消元得 \(a_k = a_{k-1}+2b_{k-2}\)。

由 \(b_{k-1} = a_{k-2}+b_{k-2}\)，得 \(b_{k-2} = b_{k-1}-a_{k-2}\)。

又 \(a_{k-1} = 2b_{k-2}\)，故 \(b_{k-2} = \frac{a_{k-1}}{2}\)，\(b_{k-1} = a_{k-2}+\frac{a_{k-1}}{2}\)。

代入 \(a_k = 2b_{k-1}\)：\(a_k = 2a_{k-2}+a_{k-1}\)，即
\[
a_k = a_{k-1} + 2a_{k-2}
\]

初始：\(a_0=1, a_1=0\)（传一次球必到邻居，无法回到出发者）。

\textbf{计算：} \(a_0=1,\ a_1=0,\ a_2=2,\ a_3=2,\ a_4=6,\ a_5=10,\ldots\)

\textbf{通项公式：} 特征方程 \(r^2 - r - 2 = 0\)，解得 \(r_1=2, r_2=-1\)。

故 \(a_k = \alpha \cdot 2^k + \beta \cdot (-1)^k\)。

由初始条件：\(\alpha+\beta=1,\ 2\alpha-\beta=0\)，解得 \(\alpha=\frac{1}{3},\ \beta=\frac{2}{3}\)。

\[
a_k = \frac{1}{3}\left[2^k + 2(-1)^k\right]
\]

\textbf{验证：} \(a_2 = \frac{1}{3}(4+2) = 2\) \(\checkmark\)，\(a_3 = \frac{1}{3}(8-2) = 2\) \(\checkmark\)。

对于一般 \(n\) 个人的情况，可类似建立含 \(n\) 个状态的递推系统，或利用对称性与特征值方法求解。

三人传球经 \(k\) 次回到出发者的方式数为 \(\boxed{\dfrac{2^k + 2(-1)^k}{3}}\)。

\begin{tipbox}[方法总结]
传球问题的关键是利用对称性简化状态：只需记录球在"出发者"和"非出发者"两个状态之间的转移。圆排列对称性大大减少了状态数。特征方程法是求解线性递推通项的标准工具。
\end{tipbox}

\vspace{6pt}

% ----- 例题7：递推计数——汉诺塔 -----
\begin{examplebox}[【例题7】递推计数——汉诺塔问题]
汉诺塔问题：有 \(n\) 个大小不同的圆盘从小到大套在柱子 A 上，每次只能移动一个圆盘，且大盘不能放在小盘上面。目标是将所有圆盘移到柱子 C 上（可借助柱子 B）。求最少移动次数 \(H_n\)。
\end{examplebox}

\textbf{解：}

\textbf{建立递推：}
\begin{enumerate}[leftmargin=1.5em]
  \item 先将上面 \(n-1\) 个圆盘从 A 移到 B（借助 C），需 \(H_{n-1}\) 次。
  \item 将最大的圆盘从 A 移到 C，需 \(1\) 次。
  \item 将 \(n-1\) 个圆盘从 B 移到 C（借助 A），需 \(H_{n-1}\) 次。
\end{enumerate}

因此 \(H_n = 2H_{n-1} + 1\)，初始条件 \(H_1 = 1\)。

\textbf{求解递推：}

改写为 \(H_n + 1 = 2(H_{n-1}+1)\)。令 \(T_n = H_n+1\)，则 \(T_n = 2T_{n-1}\)，\(T_1 = 2\)。

故 \(T_n = 2^n\)，从而
\[
H_n = 2^n - 1
\]

\textbf{验证：} \(H_1=1, H_2=3, H_3=7, H_4=15\)，均正确。汉诺塔最少移动次数为 \(\boxed{2^n-1}\)。

\begin{tipbox}[方法总结]
汉诺塔是递推的经典模型。关键是识别"最大的盘"必须单独移动，将其余盘的移动归结为相同结构的子问题。技巧：\(H_n+c = 2(H_{n-1}+c)\) 配方法求解形如 \(a_n = 2a_{n-1}+b\) 的非齐次递推。
\end{tipbox}

\vspace{6pt}

% ----- 例题8：卡特兰数——括号匹配 -----
\begin{examplebox}[【例题8】卡特兰数——括号匹配]
有 \(n\) 对括号，求所有合法的括号匹配序列的个数。例如 \(n=3\) 时，合法序列有 \(()()(),\ ()(()),\ (())(),\ (()()),\ ((()))\) 共 \(5\) 种。
\end{examplebox}

\textbf{解：}

设 \(n\) 对括号的合法序列数为 \(C_n\)。

\textbf{方法一：递推法}

考虑第一个左括号与之匹配的右括号的位置。设第一个左括号与第 \(2k+2\) 位的右括号匹配（\(k=0,1,\ldots,n-1\)），则：
\begin{itemize}[leftmargin=1.5em]
  \item 这对括号内部有 \(k\) 对括号，有 \(C_k\) 种合法序列。
  \item 这对括号外部（右侧）有 \(n-1-k\) 对括号，有 \(C_{n-1-k}\) 种合法序列。
\end{itemize}

因此
\[
C_n = \sum_{k=0}^{n-1} C_k \cdot C_{n-1-k}
\]

初始条件 \(C_0 = 1\)（空序列视为合法）。

\textbf{计算：}
\begin{align*}
C_0 &= 1 \\
C_1 &= C_0 C_0 = 1 \\
C_2 &= C_0 C_1 + C_1 C_0 = 2 \\
C_3 &= C_0 C_2 + C_1 C_1 + C_2 C_0 = 5 \\
C_4 &= C_0 C_3 + C_1 C_2 + C_2 C_1 + C_3 C_0 = 14
\end{align*}

\textbf{方法二：通项公式（反射原理）}

将左括号视为向右走一步 \((1,0)\)，右括号视为向上走一步 \((0,1)\)。合法括号序列对应从 \((0,0)\) 到 \((n,n)\) 不越过对角线的格路。

总路径数为 \(\dbinom{2n}{n}\)，越过对角线的路径数（由反射原理）为 \(\dbinom{2n}{n+1}\)。

\[
C_n = \binom{2n}{n} - \binom{2n}{n+1} = \frac{1}{n+1}\binom{2n}{n}
\]

\(n\) 对括号的合法匹配序列数为卡特兰数 \(\boxed{C_n = \dfrac{1}{n+1}\dbinom{2n}{n}}\)。

\begin{tipbox}[方法总结]
括号匹配是卡特兰数最经典的模型。递推法（"第一个左括号的匹配位置"分类）建立了卡特兰数的标准递推关系。反射原理将问题转化为格路计数，直接给出通项公式。两种方法互补，建议都掌握。
\end{tipbox}

\vspace{6pt}

% ----- 例题9：卡特兰数——出栈序列 -----
\begin{examplebox}[【例题9】卡特兰数——出栈序列]
设有一个栈，将 \(1,2,3,4,5\) 按顺序依次入栈（每次入栈后可以选择立即出栈或等待），则所有可能的出栈序列有多少种？
\end{examplebox}

\textbf{解：}

将入栈操作记为 \(+\)（向右走一步），出栈操作记为 \(-\)（向上走一步）。每次入栈必须先于对应出栈，且出栈时栈中必须有元素。这恰好对应从 \((0,0)\) 到 \((n,n)\) 不越过对角线的格路。

因此出栈序列数为卡特兰数 \(C_n\)，其中 \(n=5\)：
\[
C_5 = \frac{1}{6}\binom{10}{5} = \frac{252}{6} = \boxed{42}
\]

\textbf{验证（递推）：}
\begin{align*}
C_0&=1,\ C_1=1,\ C_2=2,\ C_3=5,\ C_4=14, \\
C_5 &= C_0 C_4 + C_1 C_3 + C_2 C_2 + C_3 C_1 + C_4 C_0 \\
&= 14 + 5 + 4 + 5 + 14 = 42 \quad \checkmark
\end{align*}

\begin{tipbox}[方法总结]
出栈序列问题是卡特兰数的另一经典应用。关键观察：入栈和出栈的操作序列可编码为括号序列（或格路），约束条件恰好对应"不越过对角线"。
\end{tipbox}

\vspace{6pt}

% ----- 例题10：卡特兰数——凸多边形三角剖分 -----
\begin{examplebox}[【例题10】卡特兰数——凸多边形三角剖分]
一个凸 \(7\) 边形，用不相交的对角线将其分割为三角形，有多少种不同的三角剖分方案？
\end{examplebox}

\textbf{解：}

凸 \(n+2\) 边形的三角剖分数等于卡特兰数 \(C_n\)。

对于凸 \(7\) 边形，\(n+2=7\)，即 \(n=5\)：
\[
C_5 = \frac{1}{6}\binom{10}{5} = \frac{252}{6} = \boxed{42}
\]

\textbf{推导（递推法）：}

设凸 \(m\) 边形的三角剖分数为 \(T_m\)。固定一条边（比如顶点 \(1\) 到顶点 \(m\)），设第三个顶点为 \(k\)（\(k=2,3,\ldots,m-1\)），则三角形 \((1,k,m)\) 将多边形分为三部分：三角形本身、左侧的 \(k\) 边形（顶点 \(1,2,\ldots,k\)）和右侧的 \(m-k+1\) 边形（顶点 \(k,k+1,\ldots,m\)）。

\[
T_m = \sum_{k=2}^{m-1} T_{k} \cdot T_{m-k+1}
\]

其中 \(T_2 = T_3 = 1\)（三角形只有一种剖分）。令 \(C_n = T_{n+2}\)，则得到标准卡特兰数递推。

\begin{tipbox}[方法总结]
凸多边形三角剖分是卡特兰数的经典几何模型。固定一条边后按"第三个顶点"分类，自然得到卡特兰数递推。注意凸 \(n+2\) 边形对应 \(C_n\)，不要搞错下标。
\end{tipbox}

\vspace{6pt}

% ----- 例题11：错排问题 -----
\begin{examplebox}[【例题11】错排问题]
有 \(5\) 封不同的信和 \(5\) 个不同的信封，每封信对应一个正确的信封。将所有信装入信封，恰好每封信都装错的装法有多少种？
\end{examplebox}

\textbf{解：}

这是经典的错排问题，需求 \(D_5\)。

\textbf{方法一：容斥原理公式}

设全集 \(U\) 为 \(5\) 封信的全部排列，\(|U|=5!=120\)。设 \(A_i\) 为第 \(i\) 封信装对的排列集合。

\begin{align*}
D_5 &= 5! - \binom{5}{1}\cdot 4! + \binom{5}{2}\cdot 3! - \binom{5}{3}\cdot 2! + \binom{5}{4}\cdot 1! - \binom{5}{5}\cdot 0! \\
&= 120 - 5\times 24 + 10\times 6 - 10\times 2 + 5\times 1 - 1\times 1 \\
&= 120 - 120 + 60 - 20 + 5 - 1 \\
&= 44
\end{align*}

\textbf{方法二：递推公式}

\begin{align*}
D_1 &= 0,\quad D_2 = 1 \\
D_3 &= 2(D_2+D_1) = 2\times 1 = 2 \\
D_4 &= 3(D_3+D_2) = 3\times 3 = 9 \\
D_5 &= 4(D_4+D_3) = 4\times 11 = 44
\end{align*}

\textbf{方法三：渐近验证}

\(D_5 \approx \dfrac{5!}{e} = \dfrac{120}{2.71828\ldots} \approx 44.15\)，与精确值 \(44\) 吻合。

恰好每封信都装错的装法有 \(\boxed{44}\) 种。

\begin{tipbox}[方法总结]
错排问题有三种计算方法：容斥原理直接求（适合较小 \(n\)）、递推公式逐步计算、渐近公式估算。考试中递推公式最实用。记住错排的标志性数值：\(D_3=2,\ D_4=9,\ D_5=44,\ D_6=265\)。
\end{tipbox}

\vspace{6pt}

% ----- 例题12：路径计数——不越对角线 -----
\begin{examplebox}[【例题12】路径计数——不越对角线的格路]
从 \((0,0)\) 出发，每次向右或向上走一步，到达 \((6,6)\)，且路径不越过对角线 \(y=x\)（即始终满足 \(y \le x\)），共有多少种走法？
\end{examplebox}

\textbf{解：}

由反射原理，从 \((0,0)\) 到 \((n,n)\) 不越过对角线的格路数为卡特兰数 \(C_n\)。

\[
C_6 = \frac{1}{7}\binom{12}{6} = \frac{1}{7} \times 924 = \boxed{132}
\]

\textbf{反射原理推导：}

总路径数为 \(\dbinom{12}{6} = 924\)。

越过对角线的路径数：由 Andr\'{e} 反射原理，等于从 \((-1,1)\) 到 \((6,6)\) 的路径数，即 \(\dbinom{12}{7} = 792\)...

等等，需要更仔细地分析。从 \((0,0)\) 到 \((6,6)\) 越过 \(y=x\) 的路径，即存在某个时刻 \(y > x\)。反射原理说：这样的路径与从 \((-1,1)\) 到 \((6,6)\) 的路径一一对应。从 \((-1,1)\) 到 \((6,6)\) 需要向右走 \(7\) 步、向上走 \(5\) 步，共 \(\dbinom{12}{7} = \dbinom{12}{5} = 792\) 条。

不越过对角线的路径数 \(= 924 - 792 = 132\)。

\textbf{验证：} \(C_6 = \dfrac{1}{7}\dbinom{12}{6} = \dfrac{924}{7} = 132\) \(\checkmark\)

\begin{tipbox}[方法总结]
"不越过对角线的格路"是卡特兰数的等价定义。反射原理的核心思想：将"非法路径"（越过对角线）通过关于 \(y=x+1\) 的反射变为从"虚起点"出发的路径，从而可直接计数。
\end{tipbox}

\vspace{6pt}

% ----- 例题13：路径计数——带障碍物 -----
\begin{examplebox}[【例题13】路径计数——带障碍物的格路]
从 \((0,0)\) 到 \((5,5)\)，每次向右或向上走一步。已知 \((2,2)\) 和 \((3,3)\) 两个点是障碍物（不能经过），且路径不越过对角线 \(y=x\)。求满足条件的路径数。
\end{examplebox}

\textbf{解：}

设 \(f(i,j)\) 为从 \((0,0)\) 到 \((i,j)\) 不越过对角线且不经过障碍物的路径数。转移方程：\(f(i,j) = f(i-1,j)+f(i,j-1)\)，在障碍物处 \(f(i,j)=0\)，且 \(j>i\) 时 \(f(i,j)=0\)（不越过对角线）。

逐步计算（行列标注，障碍物标 \(\times\)）：

\begin{center}
\renewcommand{\arraystretch}{1.3}
\begin{tabular}{c|cccccc}
  \(j\backslash i\) & 0 & 1 & 2 & 3 & 4 & 5 \\
  \hline
  0 & 1 & 1 & 1 & 1 & 1 & 1 \\
  1 & 0 & 1 & 2 & 3 & 4 & 5 \\
  2 & 0 & 0 & \(\times\) & 3 & 7 & 12 \\
  3 & 0 & 0 & 0 & \(\times\) & 7 & 19 \\
  4 & 0 & 0 & 0 & 0 & 7 & 26 \\
  5 & 0 & 0 & 0 & 0 & 0 & 26 \\
\end{tabular}
\end{center}

\textbf{关键计算说明：}
\begin{itemize}[leftmargin=1.5em]
  \item \(f(2,2)=0\)（障碍物）。
  \item \(f(3,2) = f(2,2)+f(3,1) = 0+3 = 3\)。
  \item \(f(4,2) = f(3,2)+f(4,1) = 3+4 = 7\)。
  \item \(f(3,3)=0\)（障碍物）。
  \item \(f(4,3) = f(3,3)+f(4,2) = 0+7 = 7\)。
  \item \(f(5,3) = f(4,3)+f(5,2) = 7+12 = 19\)。
  \item \(f(5,4) = f(4,4)+f(5,3) = 7+19 = 26\)。
  \item \(f(5,5) = f(4,5)+f(5,4) = 0+26 = 26\)。
\end{itemize}

满足条件的路径数为 \(\boxed{26}\)。

\begin{tipbox}[方法总结]
带障碍物的格路计数使用递推（动态规划）最为方便。逐格计算，障碍物处值为 \(0\)。注意不越过对角线的约束：当 \(j>i\) 时所有格子值为 \(0\)。多重约束叠加时，递推比公式法更可靠。
\end{tipbox}

\vspace{6pt}

% ----- 例题14：生成函数入门 -----
\begin{examplebox}[【例题14】生成函数——整数分拆与重复组合]
用生成函数方法求：将 \(n\) 个相同的球放入 \(3\) 个不同的盒子（允许空盒）的方案数。
\end{examplebox}

\textbf{解：}

设第 \(i\) 个盒子放入 \(k_i\) 个球，则 \(k_1+k_2+k_3 = n\)，\(k_i \ge 0\)。

每个盒子的生成函数为（可放入 \(0,1,2,\ldots\) 个球）：
\[
1 + x + x^2 + x^3 + \cdots = \frac{1}{1-x}
\]

三个盒子的总生成函数为（乘法原理——独立选择）：
\[
G(x) = \left(\frac{1}{1-x}\right)^3 = \frac{1}{(1-x)^3}
\]

由广义二项式展开（负二项式级数）：
\[
\frac{1}{(1-x)^3} = \sum_{n=0}^{\infty} \binom{n+2}{2} x^n = \sum_{n=0}^{\infty} \frac{(n+1)(n+2)}{2} x^n
\]

因此 \(x^n\) 的系数为：
\[
[x^n]G(x) = \binom{n+2}{2} = \frac{(n+1)(n+2)}{2}
\]

\textbf{验证（隔板法）：} 将 \(n\) 个球和 \(2\) 块隔板排列，等价于从 \(n+2\) 个位置选 \(2\) 个放隔板：\(\dbinom{n+2}{2}\)，与生成函数结果一致。

方案数为 \(\boxed{\dbinom{n+2}{2} = \dfrac{(n+1)(n+2)}{2}}\)。

\begin{tipbox}[方法总结]
生成函数方法处理"重复组合"问题的核心步骤：(1) 为每个独立变量写出其生成函数（可选值集合对应的幂级数）；(2) 相乘得到总生成函数；(3) 提取 \(x^n\) 的系数。\(\dfrac{1}{(1-x)^k}\) 的展开式 \(\displaystyle\sum \binom{n+k-1}{k-1}x^n\) 是最常用的公式。
\end{tipbox}

% ============================================================
% 第三部分：应用题
% ============================================================
\section*{\textcolor{apppurple}{三、应用题}}

\begin{appbox}[【应用1】密码学中的错排]
一个保险柜有 \(6\) 个密码转盘，每个转盘有 \(10\) 个位置（标记为 \(0\) 至 \(9\)）。初始时每个转盘恰好在正确位置。操作时，必须将每个转盘都拨到非初始位置的位置。问：有多少种操作方式？
\end{appbox}

\textbf{解：}

每个转盘有 \(9\) 个可选的非初始位置，但"每个转盘都不到初始位置"不等价于"独立选择"——因为转盘之间的选择确实是独立的。

实际上，\(6\) 个转盘互相独立，每个转盘有 \(9\) 种选择，故方案数为
\[
9^6 = \boxed{531441}
\]

\textbf{注意：} 这与错排问题不同。错排问题要求元素在不同的位置之间重新排列（一一对应），而本题每个转盘独立选择，不存在"排列"关系。

\begin{warnbox}[易错提醒]
错排问题的关键是"一一对应"：每个元素必须去某个位置，且不同元素去不同位置。如果各变量独立选择，则用乘法原理即可，不需要错排公式。
\end{warnbox}

\vspace{6pt}

\begin{appbox}[【应用2】投票问题（卡特兰数应用）]
甲、乙两人参加竞选，甲获得 \(a\) 票，乙获得 \(b\) 票（\(a > b\)）。唱票过程中，甲的累计票数始终严格大于乙的累计票数的唱票顺序有多少种？
\end{appbox}

\textbf{解：}

将每次唱甲的票记为向右走一步 \((1,0)\)，唱乙的票记为向上走一步 \((0,1)\)。唱票过程对应从 \((0,0)\) 到 \((a,b)\) 的格路。

条件"甲的累计票数始终严格大于乙"即路径始终在对角线 \(y=x\) 下方（严格），即对路径上所有中间点 \((x,y)\) 有 \(y < x\)。

由反射原理，从 \((0,0)\) 到 \((a,b)\) 越过 \(y=x\)（即存在 \(y \ge x\) 的时刻，不包括起点）的路径数等于从 \((-1,1)\) 到 \((a,b)\) 的路径数。

从 \((-1,1)\) 到 \((a,b)\) 需向右走 \(a+1\) 步、向上走 \(b-1\) 步，共 \(\dbinom{a+b}{a+1}\) 条。

所求路径数为：
\[
\binom{a+b}{a} - \binom{a+b}{a+1} = \frac{a-b}{a+b}\binom{a+b}{a}
\]

这是一个广义的卡特兰数（ballot number）。

\textbf{特例：} 若 \(a=5, b=3\)，则
\[
\frac{5-3}{5+3}\binom{8}{5} = \frac{2}{8}\times 56 = \boxed{14}
\]

这恰好等于卡特兰数 \(C_4 = 14\)（因为 \(a-b=2\) 时与 \(C_{(a-b)/2+1}\) 相关需进一步分析，但本题 \(a-b=2\) 时 \(\frac{2}{8}\times 56 = 14\) 直接得到结果）。

\vspace{6pt}

\begin{appbox}[【应用3】几何中的三角剖分]
一个凸八边形广场被若干条不相交的对角线分割为三角形区域，用于铺设不同颜色的地砖。若两种分割方案只要使用的对角线集合不同就算不同方案，则共有多少种不同的分割方案？
\end{appbox}

\textbf{解：}

凸 \(n+2\) 边形的三角剖分数为卡特兰数 \(C_n\)。凸八边形对应 \(n+2=8\)，即 \(n=6\)。

\[
C_6 = \frac{1}{7}\binom{12}{6} = \frac{924}{7} = \boxed{132}
\]

\vspace{6pt}

\begin{appbox}[【应用4】二进制序列（卡特兰数应用）]
长度为 \(10\) 的二进制序列中，有多少个序列满足：从左到右读，读到任意位置时，已读的 \(0\) 的个数不少于已读的 \(1\) 的个数？
\end{appbox}

\textbf{解：}

设序列中有 \(n\) 个 \(0\) 和 \(n\) 个 \(1\)（长度为 \(2n=10\)，即 \(n=5\)）。条件"任意位置时 \(0\) 的个数 \(\ge\) \(1\) 的个数"恰好对应合法括号序列（\(0\) 对应左括号，\(1\) 对应右括号），方案数为卡特兰数 \(C_n\)。

若序列中 \(0\) 和 \(1\) 的个数不等（比如 \(0\) 多于 \(1\)），则用 ballot number 计算。

本题要求长度为 \(10\)，\(0\) 和 \(1\) 各 \(5\) 个：
\[
C_5 = \frac{1}{6}\binom{10}{5} = \frac{252}{6} = \boxed{42}
\]

\textbf{若不要求 \(0\) 和 \(1\) 个数相等}，但要求"读到任意位置时 \(0\) 的个数 \(\ge\) \(1\) 的个数"，则需对 \(0\) 的个数 \(k \ge 5\) 分别计算 ballot number 后求和。

% ============================================================
% 第四部分：课堂练习
% ============================================================
\section*{\textcolor{tipgreen}{四、课堂练习}}

\begin{enumerate}[leftmargin=2em]
  \item 某次考试有 \(100\) 人参加，及格的有 \(70\) 人，得良的有 \(50\) 人，得优的有 \(30\) 人。已知及格且得良的有 \(35\) 人，及格且得优的有 \(20\) 人，得良且得优的有 \(15\) 人，三者都有的有 \(10\) 人。求三者都没有的人数。

  \item 求方程 \(x_1+x_2+x_3 = 15\) 满足 \(0\le x_1 \le 5,\ 0\le x_2 \le 7,\ 0\le x_3 \le 10\) 的整数解个数。

  \item 用容斥原理求 \(\varphi(210)\)。

  \item 用 \(1\times 1\) 和 \(1\times 2\) 的砖铺 \(1\times 8\) 的长条区域，共有多少种铺法？

  \item 4 个人传球，从甲开始，每次传给其他人（不能传给自己），经 \(5\) 次传球后球回到甲，有多少种传法？

  \item 求 \(C_7\)（第 \(7\) 个卡特兰数）。

  \item 凸 \(9\) 边形的三角剖分有多少种方案？

  \item \(6\) 封不同的信全部装错信封（错排），有多少种方法？

  \item 从 \((0,0)\) 到 \((7,7)\) 不越过对角线的格路有多少条？

  \item 将 \(n\) 个相同的球放入 \(4\) 个不同的盒子（允许空盒），用生成函数求方案数，并与隔板法结果对比验证。
\end{enumerate}

\vspace{10pt}

\begin{knowledgebox}[参考答案]
\textbf{1.} 三者都没有的人数：由容斥原理，
\(|A\cup B\cup C| = 70+50+30-35-20-15+10 = 90\)，故三者都没有 \(= 100-90 = \boxed{10}\) 人。

\medskip

\textbf{2.} 先算无约束解数：\(\dbinom{17}{2}=136\)。设 \(A_1\): \(x_1\ge 6\)，\(|A_1|=\dbinom{11}{2}=55\)；\(A_2\): \(x_2\ge 8\)，\(|A_2|=\dbinom{9}{2}=36\)；\(A_3\): \(x_3\ge 11\)，\(|A_3|=\dbinom{6}{2}=15\)。两两交集：\(|A_1\cap A_2|=\dbinom{5}{2}=10\)，\(|A_1\cap A_3|=1\)（\(\dbinom{2}{2}\)），\(|A_2\cap A_3|=0\)。三交集为 \(0\)。结果 \(= 136-(55+36+15)+(10+1+0)-0 = \boxed{41}\)。

\medskip

\textbf{3.} \(210=2\times 3\times 5\times 7\)。\(\varphi(210) = 210\times\frac{1}{2}\times\frac{2}{3}\times\frac{4}{5}\times\frac{6}{7} = \boxed{48}\)。

\medskip

\textbf{4.} \(a_n = a_{n-1}+a_{2}\)（斐波那契型）。\(a_1=1,\ a_2=2,\ a_3=3,\ a_4=5,\ a_5=8,\ a_6=13,\ a_7=21,\ a_8=\boxed{34}\)。

\medskip

\textbf{5.} 设 \(a_k\) 为 \(k\) 次传球后回到甲的方式数。4 人时，\(a_k = (n-1)^{k-1}\) 的修正形式。由递推 \(a_k = 3^{k-1}-a_{k-1}\cdot\)... 更精确地，\(a_k = \frac{1}{4}[3^k+3\cdot(-1)^k]\)。\(a_5 = \frac{1}{4}[243-3] = \boxed{60}\)。

\medskip

\textbf{6.} \(C_7 = \frac{1}{8}\dbinom{14}{7} = \frac{3432}{8} = \boxed{429}\)。

\medskip

\textbf{7.} 凸 \(9\) 边形对应 \(C_7 = \boxed{429}\)。

\medskip

\textbf{8.} \(D_6 = 5(D_5+D_4) = 5(44+9) = 5\times 53 = \boxed{265}\)。

\medskip

\textbf{9.} \(C_7 = \boxed{429}\) 条。

\medskip

\textbf{10.} 生成函数：\(\left(\frac{1}{1-x}\right)^4 = \frac{1}{(1-x)^4}\)，系数 \([x^n] = \dbinom{n+3}{3}\)。隔板法：将 \(n\) 个球和 \(3\) 块隔板排列，\(\dbinom{n+3}{3}\)。两者一致，方案数为 \(\boxed{\dbinom{n+3}{3} = \dfrac{(n+1)(n+2)(n+3)}{6}}\)。
\end{knowledgebox}

\vfill
\begin{center}
\small\textcolor{gray}{本讲义仅供教学参考，例题可根据学生水平调整难度。}
\end{center}

\end{document}
